\section*{Chapter 3 --- Differentiation and Convexity}

\begin{solution}{exo:tle:deriv:1}
$f'(x) = 3x^2 - 5$.

Quotient rule:
$g'(x) = \dfrac{(x^2+1) - x\cdot 2x}{(x^2+1)^2} = \dfrac{1 - x^2}{(x^2+1)^2}$.

Chain rule with $u = 2x+1$: $h'(x) = 5 \cdot 2 \cdot (2x+1)^4 = 10(2x+1)^4$.

Chain rule with $u = x^2 + 1$:
$k'(x) = \dfrac{2x}{2\sqrt{x^2+1}} = \dfrac{x}{\sqrt{x^2+1}}$.
\end{solution}

\begin{solution}{exo:tle:deriv:2}
$f'(x) = 2x - 3$, so $f(2) = -1$ and $f'(2) = 1$: the tangent is
$y = -1 + (x - 2) = x - 3$. The gap is
\[
f(x) - (x - 3) = x^2 - 4x + 4 = (x-2)^2 \geq 0 ,
\]
so the curve lies above the tangent everywhere, touching it only at $x = 2$
(as expected: $f$ is convex).
\end{solution}

\begin{solution}{exo:tle:deriv:3}
Both are difference quotients. With $f(x) = (1+x)^{100}$ at $0$:
$f'(x) = 100(1+x)^{99}$, so the limit is $f'(0) = 100$.
With $g(x) = \sqrt{x}$ at $4$: $g'(x) = \frac{1}{2\sqrt{x}}$, so the limit is
$g'(4) = \frac14$.
\end{solution}

\begin{solution}{exo:tle:deriv:4}
Domain $\R \setminus \{1\}$. Limits: at $\pm\infty$, $f(x) \sim x \to
\pm\infty$; at $1^\pm$, numerator $\to 4 > 0$ and denominator $\to 0^\pm$,
so $f \to \pm\infty$; the line $x = 1$ is a vertical asymptote.

\[
f'(x) = \frac{2x(x-1) - (x^2+3)}{(x-1)^2} = \frac{x^2 - 2x - 3}{(x-1)^2}
= \frac{(x-3)(x+1)}{(x-1)^2}.
\]
Hence $f' > 0$ on $\intoo{-\infty}{-1}$ and $\intoo{3}{+\infty}$, $f' < 0$ on
$\intoo{-1}{1}$ and $\intoo{1}{3}$. The function increases to a local maximum
$f(-1) = -2$, decreases to $-\infty$, jumps to $+\infty$ after the asymptote,
decreases to a local minimum $f(3) = 6$, then increases.
\end{solution}

\begin{solution}{exo:tle:deriv:5}
For $x \in \intoo{0}{15}$, the box has a square base of side $30 - 2x$ and
height $x$, so
\[
V(x) = x(30 - 2x)^2 .
\]
$V'(x) = (30-2x)^2 + x \cdot 2(30-2x)(-2) = (30-2x)(30 - 2x - 4x)
= (30-2x)(30-6x)$. On $\intoo{0}{15}$, $30 - 2x > 0$, so $V'$ has the sign of
$30 - 6x$: positive before $x = 5$, negative after. The volume is maximal for
$x = 5$ cm, giving $V(5) = 5 \times 20^2 = 2000\ \text{cm}^3$.
\end{solution}

\begin{solution}{exo:tle:deriv:6}
\emph{1.} $f'(x) = 4x^3 - 12x^2$ and $f''(x) = 12x^2 - 24x = 12x(x-2)$.
Thus $f'' > 0$ on $\intoo{-\infty}{0}$ and on $\intoo{2}{+\infty}$ (convex),
$f'' < 0$ on $\intoo{0}{2}$ (concave). $f''$ changes sign at $0$ and at $2$:
both are inflection points.

\emph{2.} At $a = 2$: $f(2) = 16 - 32 + 10 = -6$, $f'(2) = 32 - 48 = -16$,
tangent $y = -6 - 16(x-2) = -16x + 26$. The gap is
\[
g(x) = f(x) + 16x - 26 = x^4 - 4x^3 + 16x - 16 .
\]
Since $g(2) = g'(2) = g''(2) = 0$, $(x-2)^3$ divides $g$; division gives
$g(x) = (x-2)^3(x+2)$. Near $x = 2$, $x + 2 > 0$, so $g$ has the sign of
$(x-2)^3$: negative before $2$, positive after. The curve crosses the
tangent: $2$ is indeed an inflection point.
\end{solution}

\begin{solution}{exo:tle:deriv:7}
The function $f(x) = \sqrt{x}$ satisfies
$f''(x) = -\frac{1}{4}x^{-3/2} < 0$ on $\intoo{0}{+\infty}$: it is concave,
so its graph lies \emph{below} each tangent. The tangent at $a = 1$ is
\[
y = f(1) + f'(1)(x - 1) = 1 + \tfrac12 (x-1) = \frac{x+1}{2},
\]
whence $\sqrt{x} \leq \frac{x+1}{2}$ for all $x > 0$. Equality means the
curve touches the tangent, which happens only at the point of tangency
$x = 1$. (Equivalently: $\left(\sqrt x - 1\right)^2 \geq 0$.)
\end{solution}

\begin{solution}{exo:tle:deriv:8}
By \cref{thm:tle:deriv:convexity}, the graph of a differentiable convex
function lies above each of its tangents. The tangent at $a$ is horizontal
($f'(a) = 0$), of equation $y = f(a)$. Hence $f(x) \geq f(a)$ for all
$x \in \R$: $f(a)$ is the global minimum.
\end{solution}

\begin{solution}{exo:tle:deriv:9}
\emph{Fixed perimeter.} A rectangle with sides $x$ and $\frac{p}{2} - x$
($0 < x < \frac p2$) has area $S(x) = x\left(\frac p2 - x\right)$. Then
$S'(x) = \frac p2 - 2x$ vanishes at $x = \frac p4$, positive before and
negative after: the maximum is at $x = \frac{p}{4}$, where both sides equal
$\frac p4$ --- a square.

\emph{Fixed area.} Sides $x$ and $\frac Ax$ give perimeter
$P(x) = 2\left(x + \frac Ax\right)$, with
$P'(x) = 2\left(1 - \frac{A}{x^2}\right)$, negative for $x < \sqrt A$ and
positive after: minimum at $x = \sqrt{A}$, a square again.

\emph{Relation.} The two statements are dual. Suppose a non-square rectangle
minimized the perimeter at fixed area $A$; the square of the same perimeter
would have area $> A$ by the first statement, and shrinking it homothetically
to area $A$ would strictly decrease its perimeter --- contradicting
minimality. Each statement thus implies the other.
\end{solution}
